% A doubled-index auxiliary Pell norm and a Pell-valued fixed index.
\section{A doubled-index congruence in the Pell block}
\label{sec:doubled-index}

Two changes allow the main norm coefficient to reuse the first
exponential modulus. The auxiliary signed norm can compare doubled
Pell indices, and the final index equation can use a fixed Pell value
in place of the index itself. Together they remove both affine
parameters $A-1$ and $A+1$ from the arithmetic certificate, reducing
Theorem~\ref{thm:best101} by one addition.

Retain the fixed-four encoding of Section~\ref{sec:fixed-four} and the
common exponential witness of Section~\ref{sec:common-witness}. Write
\[
 A=a+4,\qquad D=A^2-1,\qquad J=2r+1,\qquad H_0=J+jc,
\]
where the supplied positive unknown $a$ still satisfies
\[
 a=Y_P(U+1),\qquad U=wn^2,\qquad Y_P=sn^2.
\]
The first exponential equation and its modulus remain
\begin{equation}\label{eq:doubled-first-exponent}
 M=8a+15=8A-17,\qquad d=U+ac+\gamma M.
\end{equation}

Replace the signed auxiliary norm by
\begin{equation}\label{eq:doubled-auxiliary}
 \begin{gathered}
 v=of-d^2,\qquad K=D(f^2-1),\qquad G=2K-1,\\
 v(v+1)=K(K-1)H_0^2.
 \end{gathered}
\end{equation}
The quantities $v,K,G$ are abbreviations; $v$ is a computed integer
and need not be positive. No unknown is added. If $L$ is the underlying
power-of-two exponent used to construct the coefficient index, replace
the supplied component $L$ by
\[
 T_L=\psi_4(L)
\]
and replace the final index equation by
\begin{equation}\label{eq:doubled-fixed-index}
 \kappa=T_L+\Delta a.
\end{equation}
Thus the four supplied inputs are $(x,V,H,T_L)$. The admissible index
is constructed by choosing the same finite circuit layout and exponent
$L$ as before, constructing $V,H$ from them, and supplying $\psi_4(L)$.
These choices are independent of $x$. Every other source equation is
unchanged.

\subsection{Recovering the main index from doubled coordinates}

\begin{lemma}\label{lem:doubled-main-index}
Let $A>1$, let $J>1$ be odd, and suppose $c>J$. Over positive integer
unknowns, the two norms
\[
 d^2-(A^2-1)c^2=1,
 \qquad f^2-(A^2-1)(ic^2)^2=1
\]
together with \eqref{eq:doubled-auxiliary} imply
$c=\psi_A(J)$ and $d=\chi_A(J)$. Conversely, when those two identities
hold, positive $f,i,j,o$ satisfying the auxiliary equations exist.
\end{lemma}
\begin{proof}
Pell classification gives positive indices $p,m$ such that
\[
 d=\chi_A(p),\quad c=\psi_A(p),\qquad
 f=\chi_A(m),\quad ic^2=\psi_A(m).
\]
The divisibility fact used in Proposition~\ref{prop:signed-pell-new}
gives $c\mid m$ from $c^2\mid\psi_A(m)$. The case $p=1$ is excluded
by $c>J>1$. For $p=2$ one has $\psi_A(2)=2A\ge4$; for $p\ge3$,
\[
 \psi_A(p)\ge(2A-1)^{p-1}\ge3^{p-1}\ge2p.
\]
Consequently
\begin{equation}\label{eq:doubled-comparison-bound}
 0<2p\le c\le m.
\end{equation}

Here $K>1$ and $G>1$. The new norm is exactly
\[
 (2v+1)^2-(G^2-1)H_0^2=1.
\]
Its right-hand product $K(K-1)H_0^2$ is positive, so $v\ge1$ or
$v\le-2$. Taking $I=|2v+1|$, another application of Pell
classification gives
\[
 I=\chi_G(t),\qquad H_0=\psi_G(t)
\]
for a positive index $t$. Reduction modulo $f$ gives
\[
 \begin{aligned}
 G&=2D(f^2-1)-1\equiv-\chi_A(2),\\
 2v+1&=2of-2d^2+1\equiv-\chi_A(2p).
 \end{aligned}
\]
The identities
$\chi_{-z}(t)=(-1)^t\chi_z(t)$ and
$\chi_{\chi_A(2)}(t)=\chi_A(2t)$ therefore imply
\[
 \chi_A(2t)\equiv\pm\chi_A(2p)\pmod{\chi_A(m)}.
\]
By Lemma~\ref{lem:signed-chi-new}, whose comparison hypothesis is
\eqref{eq:doubled-comparison-bound},
\[
 2t\equiv\pm2p\pmod{2m},\qquad
 t\equiv\pm p\pmod m,\qquad t\equiv\pm p\pmod c.
\]
The middle implication divides an integer equality by two; it does
not assume that two is invertible modulo $m$.

Since $c\mid f^2-1$, one has $G\equiv-1\pmod c$. The recurrence for
$\psi$ then gives
\[
 H_0\equiv\psi_{-1}(t)=(-1)^{t-1}t\pmod c.
\]
Together with $H_0=J+jc$, this proves $J\equiv\pm p\pmod c$.
Both $J$ and $p$ lie in $(0,c)$, so $J=p$ or $J+p=c$. The latter
contradicts $\psi_A(p)\equiv p\pmod2$ and the oddness of $J$.
Thus $p=J$. No oddness condition on $f$ was used in this direction.

For necessity, start with $c=\psi_A(J)$ and $d=\chi_A(J)$. Choose
\[
 m=2cJ,\qquad f=\chi_A(m),\qquad i=\frac{\psi_A(m)}{c^2}.
\]
The quotient is a positive integer: in
$(d+c\sqrt D)^{2c}$ the term with one square root contains
$(2c)c=2c^2$, and every remaining square-root term contains at least
$c^3$. Moreover $m$ is even, so
$f=2\chi_A(m/2)^2-1$ is odd. Put
\[
 K=D(f^2-1),\quad G=2K-1,\quad
 I=\chi_G(J),\quad H_0=\psi_G(J).
\]
The parameter $G$ is odd and exceeds one, so $I$ is odd and greater
than one. Because $J$ is odd,
\[
 I\equiv-\chi_A(2J)=1-2d^2\pmod f.
\]
Thus $I+2d^2-1$ is divisible by both $2$ and the odd number $f$.
The choices
\[
 v=\frac{I-1}{2}>0,\qquad
 o=\frac{I+2d^2-1}{2f}>0
\]
are integers and satisfy $v=of-d^2$. Also $G\equiv-1\pmod c$ and
$J$ is odd, whence $H_0\equiv J\pmod c$. Strict Pell growth makes
$j=(H_0-J)/c$ a positive integer. The Pell norm for $I,H_0$ gives
\eqref{eq:doubled-auxiliary}.
\end{proof}

The required hypotheses $A>1$ and $1<J<c$ are established by the
unchanged preliminary packing and interval arguments. Thus this lemma
is available at the same point as the preceding signed-block lemma,
before either exponential relation or binary decoding is used.

\subsection{Encoding the fixed index by a Pell value}

\begin{lemma}\label{lem:doubled-fixed-index}
In the new system, \eqref{eq:doubled-fixed-index} and the unchanged
second norm and gap force
$\kappa=\psi_A(L)$ and $\mu=\chi_A(L)$.
Conversely, those coordinates give a positive integer $\Delta$
satisfying \eqref{eq:doubled-fixed-index}.
\end{lemma}
\begin{proof}
After Lemma~\ref{lem:doubled-main-index}, the common-witness argument
is unchanged through the first exponential equation. The ratio lower
bound first gives $Y_P\ge U^r$ and
$a=Y_P(U+1)>U^{r+1}$. Equation~\eqref{eq:doubled-first-exponent}
then gives $U=4^J$. Since $r+1>3$, the modulus in the new index
equation satisfies the stronger bound
\begin{equation}\label{eq:doubled-modulus-growth}
 a>U^{r+1}>4^{3J}.
\end{equation}
This is a bound for $a=A-4$, not merely for $A$.

The unchanged second norm and $c=\kappa+\phi$ give
\[
 \kappa=\psi_A(t_0),\qquad \mu=\chi_A(t_0),\qquad 0<t_0<J.
\]
Since $A\equiv4\pmod a$, the integer polynomial recurrence for
$\psi$ and \eqref{eq:doubled-fixed-index} yield
\[
 \psi_4(t_0)\equiv\psi_4(L)\pmod a.
\]
The underlying encoding still satisfies $3L\le B\le n\le r$,
so $L<J$. For $1\le s<J$,
\[
 0<\psi_4(s)\le8^{s-1}<4^{3J}<a.
\]
The two residues therefore agree as integers. Strict growth of
$\psi_4$ gives $t_0=L$.

Conversely, for the intended coordinates the quotient
\[
 \Delta=\frac{\psi_A(L)-\psi_4(L)}a
\]
is integral by the same recurrence and strictly positive because
$A>4$ and $L\ge2$. Strict growth in the base follows, for example,
from the positive binomial expansion of the Pell solution.
\end{proof}

Once this index is recovered, the second exponent test proves
$q=B^L$ with the same bounds as before. The common-witness proof has
already supplied the power-of-two radix and the requisite binomial
divisibility. All fixed-four coefficient decoding is consequently
unchanged. In necessity, the new fixed component $T_L$ is chosen with
$V,H$ before querying $x$; the remaining Pell witnesses are chosen by
the common-witness construction and the two lemmas above.

\subsection{The complete count}

The earlier certificate computed the norm coefficient by
\[
 a_-=a+3,\qquad a_+=a_-+2,\qquad D=a_-a_+.
\]
Neither affine value is required after the two changes. Since
$M=8a+15$ is already computed for the first exponent test, use instead
\[
 a_2=a^2,\qquad D=a_2+M.
\]
This replaces three instructions by two. The auxiliary block has the
same length as before: use the existing $d^2$ in $v=of-d^2$;
replace the former root square and final addition of one by $v+1$
and $v(v+1)$; and use the three coefficient instructions
$K=D\,AE$, $K_-=K-1$, and $KK_-$, where the existing E16 register
is $AE=D(ic^2)^2$. The final index equation still takes one product
and one addition.

For exact comparison with the source equations, let
\[
 F_{16}=f^2-AE-1,\qquad
 K_s=D(f^2-1),\qquad K_c=D\,AE.
\]
The changed norm has the triangular identity
\[
 F_{17}^{\mathrm{calc}}
 =F_{17}^{\mathrm{source}}
  +D F_{16}(K_s+K_c-1)H_0^2.
\]
The executable certificate checks this identity, the inherited E7 and
E18 corrections from $\theta=B-4$, every other residual, and every
primitive instruction. The fixed-index and positive-domain conclusions
are supplied by the proofs above, independently of these symbolic
checks.

\begin{theorem}\label{thm:best100-pell}
For every represented set, an admissible fixed index $(V,H,T_L)$
can be chosen such that the system described above represents it over
34 positive integer unknowns with 22 equality tests. Its straight-line
certificate has
\[
 \boxed{100=56\text{ multiplications}+44\text{ additions}}.
\]
Fixed numerals and supplied index parameters are free. Relative to
Theorem~\ref{thm:best101}, the saving is one addition.
\end{theorem}

The full source equations, primitive schedule, domains and residual
identities are recorded by
\texttt{round24\_1980\_doubled\_pell\_certificate.py} and its JSON receipt.
